\subsection{Uniform completion of the finite-alteration exponent set}

\begin{statement}{A.71}{All finite alterations of all integers on every elliptic cubic twist}{Proved}
Let $E_d$, $O$, $R$ and $\Gamma_3$ be as in A.70.
For every integer $n$ and every $a\in\widehat\Z$
equal to $n$ outside finitely many coefficient primes,
$a\kappa(R)$ has a $\Gamma_3$ Selmer lift if and only
if $a=n$.

Thus the entire finite-alteration classification of
A.69 holds uniformly on every rational-origin cubic
twist carrying a nontorsion rational point, including
arbitrary finite support around zero. In this family
the first listed finite congruence refinement already
excludes every noninteger exponent. Infinite
alterations, the quotient's fibres, and cubic torsors
without a rational origin are not included.
\end{statement}
\begin{proof}
A.70 handles every nonzero integer baseline.
It remains to exclude nonzero finite-support
exponents. To avoid confusion with the exponent
$a$, write $u^3=d$. In the short model put
\[
 E_d:Y^2=X^3-432d^2,\qquad K=\Q(\zeta),
 \qquad L=K(u)=\Q(E_d[3]).
\]
Indeed the three-division polynomial is
$3X(X^3-1728d^2)$: its nonzero roots give
the points $(12u^2\zeta^i,\pm36d)$,
and $X=0$ gives $Y=\pm12d\sqrt{-3}$.
The normalized retained flex ratio is
\[
 g_d=\frac{6u\zeta^2X-Y-36d}{6uX-Y-36d}.
 \tag{A71.1}
\]
Its tangent-line divisor is $3(d_1)-3(d_0)$,
where $d_i=(12u^2\zeta^i,36d)$, and
$g_d(O)=1$. Thus it is exactly the actual
$\Gamma_3$ character of A.67.

We apply the common joint-tower argument
(A69.2)--(A69.6) to this curve and this
unit. That argument uses CM, a nontorsion
rational point, the normalized retained
character, a finite abelian base $F_0/K$,
and a torsion value $h\in F_0^\times$
whose quadratic radical is not Galois
over $K$. Odd translation kernels kill
quadratic characters automatically; when
two occurs in the support, a residual
$C_3$ supplies (A69.3). No special
equation of $E_9$ was used in that
common argument.

\par\smallskip\noindent\textbf{Support containing coefficient two.}\enspace
Let $c^3=2$ and $F_0=K(c)$. The
nonzero two-torsion point
\[
 S_2=(6cu^2,0)
\]
satisfies the exact, twist-independent
identity
\[
 g_d(S_2)=\frac{\zeta^2c-1}{c-1}=h.
 \tag{A71.2}
\]
This is the same nonsquare at $31$
as in (A69.7), with three cubic
conjugates of product one. Its
quadratic radical is therefore not
Galois over $K$.
Also
\[
 F_0(E_d[2])=F_0(u)=F_0L:
\]
the two-torsion $X$-coordinate gives
$u^2$, which determines $u$ because
$u^3=d$. The residual image is
$C_3$ unless $u\in F_0$.
That exception occurs exactly when
the rational cube class of $d$
is one of $1,2,4$. Kummer theory
over $K$ identifies the one cyclic
cubic field $K(c)$, and a rational
number becoming a cube over the
quadratic extension $K/\Q$ was
already a rational cube.

These three exceptional twists
cannot contain the stipulated
nontorsion $R$. The certified
rows of A.38 give respectively
conductors $27,36,108$ and
nonzero central-value enclosures
\[
\begin{array}{c|l}
d&L(E_d,1)\\\hline
1&[0.588879583428483319104563\ \pm\ 2.77\cdot10^{-25}]\\
2&[0.701091052662727130587510\ \pm\ 5.35\cdot10^{-25}]\\
4&[1.1129126745223053845620\ \pm\ 3.84\cdot10^{-23}].
\end{array}
\]
Their Mordell--Weil ranks are zero
by the CM theorem used in A.37
\citep[Corollary C and Remark (3)]{AMRubin1987}.
Consequently the required residual
$C_3$ survives in every case at
hand, and the common tower
argument excludes all finite
support containing two.

\par\smallskip\noindent\textbf{Odd support containing a prime $p\ge5$.}\enspace
Take the fixed witness on $E_9$
provided uniformly in $p$ by A.69:
\[
 b^3=3,\qquad M_0=K(b)(E_9[p]),
 \qquad S_p\in E_9[p],\qquad h=g_0(S_p).
\]
The element $h$ is nonsquare in
$M_0$. There is an automorphism
$\tau$ fixing $E_9[p]$ and
cycling $b$, and
$h\tau(h)\tau^2(h)=1$.
Now set
\[
 F_0=M_0K(u).
 \tag{A71.3}
\]
It is abelian over $K$, and
$[F_0:M_0]$ is one or three,
so $h$ remains nonsquare.
Choose any lift $\sigma$ of
$\tau$ to $\operatorname{Gal}(F_0/K)$.
The lift need not have order
three. Since $h\in M_0$, however,
\[
 h\sigma(h)\sigma^2(h)=1.
\]
If $F_0(\sqrt h)/K$ were Galois,
the squareclass of $h$ would be
$\sigma$-invariant. The last
identity would make $[h]^3=[h]$
trivial, a contradiction. Thus
the required non-Galois radical
persists without any disjointness
assumption on $K(u)$ and $M_0$.

The element $t=b^2/u\in F_0$
satisfies $t^3=9/d$. The scaling
\[
 (X,Y)\longmapsto(t^2X,t^3Y)
\]
identifies $E_d$ with the short
model of $E_9$ and transports
$g_d$ to its normalized flex
ratio. Under the further exact
change $X=36x$, $Y=216y+108$,
this ratio is the $g_0$ of A.69.
The inverse image $S_d$ of $S_p$
is consequently a $p$-torsion
point defined over $F_0$ with
$g_d(S_d)=h$. It is interior
because $p\ge5$, and the base
already contains $L$.
Since the support is odd,
all translation kernels kill
quadratic characters without a
residual-two condition. The
common tower argument applies
and excludes this case as well.

The only nonempty support not
yet considered is $\{3\}$,
excluded uniformly by A.70.
Zero is realized by the
rational $O$-cusp, and every
nonzero integer by its rational
point. This proves both directions.

The exact flex-line, two-torsion,
and transport identities are
checked by the executed program
\texttt{code/a\_main\_twist\_transport.py}.
Its output also identifies the
three previously certified
rank-zero rows; it does not
replace their A.38 enclosures
by a numerical-rank heuristic.
\end{proof}
